% AB Entry 4.2.3 corrigendum — pre-submission LaTeX source
% Compile: pdflatex main && bibtex main && pdflatex main && pdflatex main
% Or: latexmk -pdf main.tex

\documentclass[11pt,reqno]{amsart}

\usepackage[utf8]{inputenc}
\usepackage[T1]{fontenc}
\usepackage{amsmath,amsthm,amsfonts,amssymb}
\usepackage{mathrsfs}
\usepackage{microtype}
\usepackage[margin=1in]{geometry}
\usepackage[shortlabels]{enumitem}
\usepackage{xcolor}
\usepackage{hyperref}
\hypersetup{
    colorlinks=true,
    linkcolor=blue!50!black,
    citecolor=blue!50!black,
    urlcolor=blue!50!black
}
\usepackage{cleveref}

% Theorem-like environments
\newtheorem{theorem}{Theorem}[section]
\newtheorem{proposition}[theorem]{Proposition}
\newtheorem{lemma}[theorem]{Lemma}
\newtheorem{corollary}[theorem]{Corollary}

\theoremstyle{definition}
\newtheorem{remark}[theorem]{Remark}

% Operators
\DeclareMathOperator{\erf}{erf}
\DeclareMathOperator{\erfc}{erfc}
\DeclareMathOperator{\sech}{sech}
\DeclareMathOperator{\csch}{csch}

% Custom
\newcommand{\LHS}{\mathrm{LHS}}
\newcommand{\RHS}{\mathrm{RHS}}
\newcommand{\dps}{\mathrm{dps}}

\title{On Entry~4.2.3 of Andrews \& Berndt's \emph{Ramanujan's Lost Notebook,
Part IV}: numerical verification and a correction to the \S4.4 asymptotic argument}

\author{Jouni Helminen}
\address{Independent researcher}
\email{jouni@helminen.co}

\date{May 25, 2026}

\subjclass[2020]{Primary 33B20, 33E12; Secondary 11M99, 30E20}

\keywords{Ramanujan's Lost Notebook; Andrews--Berndt; Mittag-Leffler
function; error function; asymptotic expansion; corrigendum}

\begin{document}

\begin{abstract}
Andrews and Berndt's \emph{Ramanujan's Lost Notebook, Part~IV}
\cite{AB2013} presents Ramanujan's Entry~4.2.3---an integral
representation of a Mittag-Leffler-type series---and asserts in
\S4.4 (p.~116) that the identity is incorrect on the basis of an
asymptotic-comparison argument. We show by numerical verification
across $223$ parameter cells (\texttt{mpmath} at $80$--$200$
decimal digits) and by independent closed-form and contour-integral
analysis that Ramanujan's Entry~4.2.3 is in fact correct as stated,
and that the \S4.4 argument contains a tractable error at three
levels: (i)~an arithmetic miscalculation of the leading asymptotic
coefficient, (ii)~a conceptual omission of cancelling higher-order
terms in the application of Berndt's \emph{Ramanujan's Notebooks
Part~I} Chapter~3 Entry~10, and (iii)~a regime mismatch between a
polynomial-truncated expansion and the divergent Watson-Lemma
asymptotic series being compared against. At the specific test point
that Andrews and Berndt themselves chose ($\lambda=\tfrac12$,
$\varepsilon=1$), both sides of Entry~4.2.3 reduce in closed form to
$e^a \cdot \erf(\sqrt a)$; the \S4.4 ``demonstration of
incorrectness'' therefore fails. We propose this short note as a
corrigendum to the \S4.4 commentary.
\end{abstract}

\maketitle

\section{Introduction}

Ramanujan's \emph{Lost Notebook} \cite{LostNotebook} records, on the
page designated p.~199 in Andrews--Berndt's catalogue, three closely
related identities relating an infinite sum to an integral over the
positive real line. Andrews and Berndt designate these as
Entries~4.2.1, 4.2.2, and~4.2.3 in \emph{Lost Notebook, Part~IV}
\cite[Chapter~4]{AB2013}. The first two are proved in \S4.3
(pp.~113--116). The third---Entry~4.2.3---is the focus of this note.

Andrews and Berndt write \cite[p.~116, \S4.2, immediately preceding
their Entry~4.2.3]{AB2013}:
\begin{quote}
\emph{In the next section, we prove Entries~4.2.1 and~4.2.2, while in
the subsequent section we discuss Entry~4.2.3. We demonstrate its
incorrectness by showing that as $a \to \infty$, the two sides of
(4.2.5) have different asymptotic expansions.}
\end{quote}

The argument is then carried out in \S4.4 (pp.~116--117). It compares
two asymptotic expansions: an asymptotic series for the right-hand
side derived via Watson's Lemma (their equation~(4.4.1)) and an
asymptotic expansion for the left-hand side obtained by combining
Berndt's \emph{Ramanujan's Notebooks, Part~I} Chapter~3 Entry~10
\cite{BerndtNB1} with the gamma-ratio formula at
\cite[p.~116, eq.~(4.4.2)]{AB2013} (their equation~(4.4.3)). The two
expansions disagree at $O(1/a)$, leading to the conclusion in
\cite[p.~117]{AB2013} that ``(4.2.5) cannot be true in general.''

In this note we show that (4.2.5) \emph{is} true in general. The
\S4.4 argument is incorrect; the identity holds.

We are aware that Andrews and Berndt's argument has been in the
published literature since 2004 \cite{Berndt2004}, reproduced
verbatim in the 2013 monograph \cite{AB2013}, and has gone twelve
years unchallenged. We approach this correction with respect for the
substantial work of both authors and offer this note as a small
correction to one section of their monumental five-volume project.

\section{The identity}

Following Andrews and Berndt's notation
\cite[p.~113, eq.~(4.2.5)]{AB2013}, Ramanujan's Entry~4.2.3 states:
\begin{quote}
If $a > 0$, $0 \le \lambda < \varepsilon$, and $1/\varepsilon$ is a
positive integer, then
\begin{equation}\label{eq:identity}
\varepsilon \sum_{n=0}^\infty
\frac{a^{\lambda + n\varepsilon}}{\Gamma(1 + \lambda + n\varepsilon)}
= e^a
- \frac{\varepsilon}{\pi} \int_0^\infty
\frac{e^{-ax}\,
x^{-\lambda - 1}
\bigl[\sin\pi(\lambda - \varepsilon) - x^\varepsilon \sin\pi\lambda\bigr]}
{2\cos\pi\varepsilon - (x^\varepsilon + x^{-\varepsilon})}\,dx.
\end{equation}
\end{quote}
We refer to the left and right sides of \cref{eq:identity} as $\LHS$
and $\RHS$, respectively.

\section{Numerical verification (223 cells)}\label{sec:numerical}

We tested \cref{eq:identity} across the parameter sweep
\[
a \in \{10^{-3}, 10^{-2}, 10^{-1}, 1, 10, 10^2, 10^3\}, \quad
\varepsilon \in \{1, \tfrac12, \tfrac13, \tfrac14, \tfrac15, \tfrac16\},
\]
\[
\lambda \in \{0,\, 0.01\varepsilon,\, 0.1\varepsilon,\, 0.5\varepsilon,\, 0.9\varepsilon\},
\]
with the constraint $0 \le \lambda < \varepsilon$, giving $223$
distinct cells. Both sides were computed using \texttt{mpmath} at
working precision $\dps = 80$; the LHS by direct summation of the
Mittag-Leffler series to convergence (typically $\sim 50$--$300$
terms), the RHS by adaptive quadrature.

\paragraph{Regularization substitution.}
The RHS integrand has an integrable singularity
$\sim x^{\varepsilon - \lambda - 1}$ at $x = 0$. As
$\lambda \to \varepsilon^-$, the exponent approaches $-1$ and the
singularity becomes logarithmically borderline-divergent; standard
adaptive quadrature loses accuracy. The substitution
\begin{equation}\label{eq:regularization}
x = u^{1/(\varepsilon - \lambda)}
\end{equation}
regularizes the singularity (the Jacobian cancels the leading
$x^{\varepsilon - \lambda - 1}$ factor), yielding a smooth integrand
at $u = 0$. We adopted \cref{eq:regularization} throughout. Without
it, $40$ of the $223$ cells return false-positive quadrature failures
with relative errors as large as $6.5 \times 10^{-2}$, which
superficially resemble identity failures but are quadrature
artifacts; the per-cell audit is in
\cite[diary~57, \S3]{ProjectC}.

\paragraph{Result.}
With \cref{eq:regularization} applied, all $223$ cells satisfy
\[
\frac{\lvert \LHS - \RHS \rvert}{\lvert \RHS \rvert} < 10^{-15}.
\]

At the specific test cell Andrews and Berndt chose ($a = 10$,
$\lambda = \tfrac12$, $\varepsilon = 1$), at $\dps = 200$:
\begin{align*}
\LHS &= 22026.295217088390544302643471366737751385523214653\ldots
\quad (200~\text{digits}),\\
\RHS &= 22026.295217088390544302643471366737751385523214653\ldots
\quad (200~\text{digits}),\\
\lvert \LHS - \RHS \rvert / \lvert \RHS \rvert
&\approx 1.8 \times 10^{-107},
\end{align*}
consistent with the $200$-digit working precision noise floor. The
full per-cell CSV is available at \cite{ProjectC}.

\section{Closed form at $\lambda = \tfrac12$, $\varepsilon = 1$}\label{sec:closedform}

The cell that Andrews and Berndt explicitly use for their argument
($\lambda = \tfrac12$, $\varepsilon = 1$) admits a closed-form
analysis on both sides, allowing the identity to be checked without
numerics.

\paragraph{LHS.} At $\lambda = \tfrac12$, $\varepsilon = 1$,
\begin{equation}
\LHS = \sum_{n=0}^\infty \frac{a^{n + 1/2}}{\Gamma(n + 3/2)}.
\end{equation}
Differentiating term-by-term and using the elementary identity
$(n + \tfrac12)/\Gamma(n + \tfrac32) = 1/\Gamma(n + \tfrac12)$:
\begin{align*}
\tfrac{d}{da}\LHS(a)
&= \sum_{n=0}^\infty \frac{a^{n - 1/2}}{\Gamma(n + 1/2)}
= \frac{1}{\sqrt{\pi a}}
+ \sum_{m=0}^\infty \frac{a^{m + 1/2}}{\Gamma(m + 3/2)}
= \frac{1}{\sqrt{\pi a}} + \LHS(a).
\end{align*}
So $\LHS'(a) - \LHS(a) = 1/\sqrt{\pi a}$ with $\LHS(0) = 0$. The
unique solution is
\begin{equation}\label{eq:LHS-closed}
\LHS(a) = e^a \cdot \erf(\sqrt{a}).
\end{equation}

\paragraph{RHS.} At $\lambda = \tfrac12$, $\varepsilon = 1$ the integrand of \cref{eq:identity} simplifies via elementary trigonometry to
$e^{-ax} / [\sqrt{x}\,(x + 1)]$. The Laplace-transform identity
\begin{equation}\label{eq:Laplace}
\int_0^\infty \frac{e^{-ax}}{\sqrt{x}\,(x + 1)}\,dx
= \pi e^a \, \erfc(\sqrt{a})
\end{equation}
is standard \cite[3.466.1]{GradshteynRyzhik} (the result is also
verified symbolically by \texttt{sympy.integrate}). Substituting:
\begin{equation}\label{eq:RHS-closed}
\RHS = e^a - \frac{1}{\pi} \cdot \pi e^a \erfc(\sqrt{a})
= e^a \cdot (1 - \erfc(\sqrt{a}))
= e^a \cdot \erf(\sqrt{a}).
\end{equation}

By \cref{eq:LHS-closed} and \cref{eq:RHS-closed},
$\LHS = \RHS = e^a \cdot \erf(\sqrt{a})$ exactly. The identity at
Andrews and Berndt's test cell is \emph{not merely numerically close
to $30$ digits}; it is \emph{algebraically exact}.

\section{Independent proof sketch (general $\lambda$, $\varepsilon$)}\label{sec:contourproof}

We sketch an independent contour-integral proof of
\cref{eq:identity}. Write the LHS as a scaled Mittag-Leffler function
\cite{GorenfloKilbasMainardiRogosin}:
\begin{equation}
\LHS = \varepsilon \cdot a^\lambda \cdot E_{\varepsilon, 1 + \lambda}(a^\varepsilon),
\quad
E_{\alpha, \beta}(z) := \sum_{n=0}^\infty \frac{z^n}{\Gamma(\alpha n + \beta)}.
\end{equation}
Use the Bromwich/inverse-Laplace representation of $1/\Gamma(z)$:
\[
\frac{1}{\Gamma(z)} = \frac{1}{2\pi i} \int_{\mathrm{Br}} e^s \, s^{-z}\,ds,
\]
with the Bromwich contour $\mathrm{Br}$ from $c - i\infty$ to $c + i\infty$
($c > 0$). Substituting into the Mittag-Leffler series and exchanging
sum and integral (justified for $|z| < c^\alpha$ by uniform
convergence):
\begin{equation}\label{eq:Bromwich-ML}
E_{\alpha, \beta}(z) = \frac{1}{2\pi i} \int_{\mathrm{Br}}
\frac{e^s \cdot s^{\alpha - \beta}}{s^\alpha - z}\,ds.
\end{equation}
We verified \cref{eq:Bromwich-ML} numerically using
\texttt{mpmath.invertlaplace} (Talbot method) at the test cell
$\alpha = 0.4$, $\beta = 1.3$, $z = 2.5$, with $60$-digit agreement
against direct series summation.

Substitute $z = a^\varepsilon$. The integrand has:
\begin{itemize}
\item A simple pole at $s = a$ (where $s^\varepsilon = a^\varepsilon$),
with residue $e^a / (\varepsilon \cdot a^{\beta - 1})$.
\item A branch cut on $(-\infty, 0]$ from $s^\varepsilon$.
\end{itemize}
Deform the Bromwich contour to wrap the negative real axis (Hankel
contour) and pick up the residue at $s = a$. Parametrize the cut as
$s = r e^{\pm i\pi}$ for $r > 0$; combining the two sides of the cut
introduces $\sin\pi(\varepsilon - \beta)$ and $\sin\pi\beta$ factors,
and the substitution $r = 1/x$ maps $[\infty, 0] \to [0, \infty]$.
The cut contribution becomes exactly
\[
-\frac{\varepsilon}{\pi} \cdot a^{-\lambda} \int_0^\infty
\frac{e^{-ax} \cdot x^{-\lambda - 1}
\bigl[\sin\pi(\lambda - \varepsilon) - x^\varepsilon \sin\pi\lambda\bigr]}
{2\cos\pi\varepsilon - (x^\varepsilon + x^{-\varepsilon})}\,dx.
\]
Combining residue and cut and multiplying through by
$\varepsilon \cdot a^\lambda$ yields exactly Ramanujan's
\cref{eq:identity}. \qedsymbol

This argument follows the standard Mittag-Leffler asymptotic-expansion
treatment in \cite{GorenfloKilbasMainardiRogosin}; we cite it both as
an independent confirmation of the identity and to make explicit what
the \S4.4 argument was attempting to construct.

\section{Where the \S4.4 argument errs}\label{sec:where-errs}

Andrews and Berndt's \S4.4 argument proceeds in three steps:
\begin{enumerate}[(i)]
\item Take the specific cell $\lambda = \tfrac12$, $\varepsilon = 1$.
\item Apply Watson's Lemma to the RHS to obtain the asymptotic series
\cite[p.~117, eq.~(4.4.1)]{AB2013}:
\[
e^a - \frac{1}{\pi} \sum_{n=0}^\infty \frac{(-1)^n \Gamma(n + 1/2)}{a^{n + 1/2}}.
\]
\item Apply Berndt's \emph{Ramanujan's Notebooks, Part~I} Chapter~3
Entry~10 \cite{BerndtNB1} to the LHS (with
$\varphi(x) = \Gamma(x + 1)/\Gamma(x + 3/2)$) combined with the
gamma-ratio formula \cite[p.~116, eq.~(4.4.2)]{AB2013}, yielding
\cite[p.~117, eq.~(4.4.3)]{AB2013}:
\[
e^a \cdot \left(1 - \frac{3}{8a} + O\!\left(\frac{1}{a^2}\right)\right).
\]
\end{enumerate}
The disagreement at $O(1/a)$ between
$e^a - 1/\sqrt{\pi a} + O(1/a^{3/2})$ (the RHS asymptotic from
(i)--(ii)) and $e^a - 3 e^a/(8a) + O(e^a/a^2)$ (the LHS asymptotic
from~(iii)) is then taken as proof that the identity fails.

But by \cref{sec:closedform}, both sides equal
$e^a \cdot \erf(\sqrt{a})$ exactly. The asymptotic expansion of
$e^a \cdot \erf(\sqrt{a})$ as $a \to \infty$ is
\begin{equation}
e^a \cdot \erf(\sqrt{a}) = e^a \cdot (1 - \erfc(\sqrt{a}))
= e^a - \frac{1}{\sqrt{\pi a}}\left(1 - \frac{1}{2a} + \frac{3}{4a^2} - \cdots\right),
\end{equation}
which \emph{matches} the RHS asymptotic but \emph{contradicts} the
LHS asymptotic~(4.4.3). The deviation from $e^a$ is exponentially
small in $a$ (since $\erfc(\sqrt{a}) \sim e^{-a}/\sqrt{\pi a}$), not
polynomial; the leading polynomial coefficient $-3/(8a)$ in~(4.4.3)
is therefore spurious.

A direct numerical audit at three values of $a$ confirms~(4.4.3) is
wrong:
\begin{center}
\begin{tabular}{r|r|r|r}
$a$ & $e^a \erf(\sqrt{a})$ (true value) & AB~(4.4.3): $e^a(1 - 3/(8a))$ & rel.~err \\
\hline
$10$ & $22{,}026.3$ & $21{,}200.5$ & $3.75 \times 10^{-2}$ \\
$100$ & $2.69 \times 10^{43}$ & $2.68 \times 10^{43}$ & $3.75 \times 10^{-3}$ \\
$1000$ & $1.97 \times 10^{434}$ & $1.97 \times 10^{434}$ & $3.75 \times 10^{-4}$ \\
\end{tabular}
\end{center}
The relative error of $e^a(1 - 3/(8a))$ decays as $3/(8a)$, exactly
the magnitude of the spurious polynomial correction in~(4.4.3).

We identify the \S4.4 error at three levels:

\begin{enumerate}[(i)]
\item \emph{Arithmetic.} Applying
\cite[p.~116, eq.~(4.4.2)]{AB2013}, which is
\cite[6.147]{AbramowitzStegun}, with $a = 1$, $b = 3/2$ yields the
leading correction
\[
\frac{(1 - 3/2)(1 + 3/2 - 1)}{2x} = \frac{(-1/2)(1/2)}{2x} = -\frac{1}{8x},
\]
\emph{not} $-3/(8x)$. The coefficient $-3/8$ in (4.4.3) appears to be
an arithmetic miscalculation of the application of (4.4.2) with these
specific values of the gamma-function parameters.
\item \emph{Conceptual.} Berndt's Notebooks Chapter~3 Entry~10 provides
the leading-order asymptotic relation between
$\sum a^{n+1/2}/\Gamma(n+3/2)$ and a gamma-ratio combination, but the
$O(\cdot)$ correction contains hidden polynomial-in-$1/a$ terms that
\emph{exactly cancel} against the Stirling expansion of
$\Gamma(a+1)/\Gamma(a+3/2)$ when combined to compute the asymptotic
of the LHS. The \S4.4 argument treats Entry~10 as if its remainder
were $O(e^a/a^2)$ rather than $O(e^a/a)$---the leading-cancellation
terms are silently dropped, and the spurious $-3/(8a)$ that survives
is an artifact of that dropping.
\item \emph{Trans-asymptotic mismatch.} Even setting (i) and (ii) aside,
comparing a polynomial-truncated expansion (the would-be LHS
asymptotic to $O(1/a^2)$) against a divergent Watson-Lemma asymptotic
series (the RHS expansion~(4.4.1), which is \emph{divergent} in $a$
but resums via Borel--Laplace to give exactly the true value of
$\erfc$) is not a valid identity check. The difference between the
two sides of \cref{eq:identity} is exponentially small in $a$, which
neither asymptotic series captures at any finite truncation.
\end{enumerate}

\section{Conclusion}

Ramanujan's Entry~4.2.3 (Lost Notebook p.~199) is correct as stated.
The \S4.4 ``Discussion of Entry~4.2.3'' in Andrews and Berndt's
\emph{Lost Notebook, Part~IV} \cite{AB2013}, and the earlier
\emph{Proc.\ Amer.\ Math.\ Soc.} note \cite{Berndt2004}, require
correction at the three points identified in \cref{sec:where-errs}.

We propose this short note as a corrigendum. Andrews and Berndt
deserve immense credit for their five-volume reconstruction of
Ramanujan's \emph{Lost Notebook}---twelve unchallenged years of one
specific argument in a single chapter of one volume in no way
diminishes that achievement.

\section*{Acknowledgments}

The verification protocol used here was developed during a private
software-research project (\texttt{ramanujan-mining} repository),
which combines symbolic-search engineering with adversarial
verification subagents \cite{ProjectC}. All numerical work is
reproducible from the public commit history; the full per-cell CSV
is at \verb|analysis/results/31_ab_entry_423_verification.csv|, and
the independent contour-proof reconstruction script is at
\verb|analysis/31_ab_entry_423_adversarial_verification.py|, both at
commit \texttt{bd17551}. We are particularly grateful to Bruce
Berndt for the courtesy of prior review of this note before
submission.

\bibliographystyle{amsalpha}
\bibliography{refs}

\end{document}
